% ============================================
% Chapter 4: Data Exploration and Preparation
% ============================================
% Rewritten — June 2026
% Based on Khadija's ML pipeline work, retrained models

\chapter{Data Exploration and Preparation}

\section{Introduction}

Before any machine learning model can be trained, the data must be understood, cleaned, and properly prepared. This chapter walks through the complete data preparation pipeline that transforms raw network flow records from the CICIDS2017 dataset into a clean, balanced, and standardized dataset ready for model training. Every decision made at this stage --- how to handle missing values, what to do with outliers, whether to balance the classes --- directly impacts the quality of the detection system. For this reason, we follow the PDCA (\textit{Plan-Do-Check-Act}) methodology, documenting not only what was done but why each choice was made.

The primary objective is to produce a dataset that enables the detection pipeline to meet the success criteria defined in the project plan:
\begin{itemize}
    \item \textbf{F1-Score} $\geq$ 88\%
    \item \textbf{Accuracy} $\geq$ 90\%
    \item \textbf{False Positive Rate (FPR)} $<$ 6\%
\end{itemize}

\section{Dataset Overview}

\subsection{What is CICIDS2017?}

The \textbf{CICIDS2017} dataset (\textit{Canadian Institute for Cybersecurity Intrusion Detection System}, 2017) is one of the most widely cited academic benchmarks for network intrusion detection research. Unlike older datasets such as KDD-99 --- which dates back to 1999 and no longer reflects modern network behavior --- CICIDS2017 was generated in a controlled environment that replicates realistic network traffic conditions, including both benign user activity and simulated attacks across multiple days.

\begin{table}[H]
\centering
\caption{Dataset characteristics — CICIDS2017 PortScan subset}
\label{tab:dataset}
\begin{tabular}{ll}
\toprule
\textbf{Property} & \textbf{Value} \\
\midrule
Number of flow records (rows) & 286,467 \\
Number of raw features (columns) & 79 \\
\texttt{PortScan} class & 158,930 (55.48\%) \\
\texttt{BENIGN} class & 127,537 (44.52\%) \\
File used & \texttt{Friday-WorkingHours-Afternoon-PortScan.pcap\_ISCX.csv} \\
Format & CSV — aggregated network flows (\textit{flow-level}) \\
\bottomrule
\end{tabular}
\end{table}

\subsection{Understanding the Class Distribution}

The dataset contains 55.48\% PortScan traffic and 44.52\% benign traffic. This is a slight imbalance in favor of the attack class, which is unusual: in most intrusion detection datasets, attacks represent a small minority (typically $<$ 5\%). Here, the PortScan class actually constitutes the \textit{majority}. This has important consequences:

\begin{itemize}
    \item \textbf{For supervised models} (Random Forest, XGBoost): the imbalance is moderate and can be managed with SMOTE (Section~\ref{sec:smote}).
    \item \textbf{For unsupervised models} (Isolation Forest): the majority class being the attack class creates a fundamental problem, which we analyse in detail in Chapter~5 (\textit{majority anomaly paradox}).
\end{itemize}

\begin{figure}[H]
\centering
\includegraphics[width=0.65\textwidth]{class_distribution.png}
\caption{Class distribution in the CICIDS2017 PortScan dataset. The attack class (PortScan) constitutes 55.48\% of the total records --- a slight majority.}
\label{fig:class_dist}
\end{figure}

\section{Feature Selection}

\subsection{From the Data Dictionary to the CSV}

The project's Data Dictionary defines 13 features theoretically relevant for detecting port scan attacks. However, CICIDS2017 is a collection of aggregated network flows (\textit{flow-level data}), not raw packet captures. As a result, not all 13 features are available from the CSV. A mapping process was conducted to find the best proxy for each theoretical feature.

\begin{longtable}{p{3.5cm}cp{3.5cm}p{4.5cm}}
\caption{Feature mapping: Data Dictionary $\rightarrow$ CICIDS2017 CSV columns} \\
\label{tab:feature_mapping} \\
\toprule
\textbf{Feature (Dictionary)} & \textbf{Avail.} & \textbf{CSV Column} & \textbf{Justification} \\
\midrule
\endfirsthead
\toprule
\textbf{Feature (Dictionary)} & \textbf{Avail.} & \textbf{CSV Column} & \textbf{Justification} \\
\midrule
\endhead
Distinct Dst Ports & $\approx$ & Destination Port & Proxy: scanners target many different ports, so the port number itself carries discriminative signal \\
SYN Flag Count & \checkmark & SYN Flag Count & Direct match — counts SYN packets per flow \\
RST Flag Count & \checkmark & RST Flag Count & Direct match — counts RST (reset) packets per flow \\
Flow Duration & \checkmark & Flow Duration & Direct match — scan flows are typically very short \\
Total Fwd Packets & \checkmark & Total Fwd Packets & Direct match — scanners send few packets per probe \\
ACK Flag Count & \checkmark & ACK Flag Count & Direct match — counts ACK packets per flow \\
IAT Mean & \checkmark & Flow IAT Mean & Direct match — inter-arrival time between packets \\
Bwd Packet Length & \checkmark & Bwd Packet Length Mean & Mean variant — average backward packet size \\
TCP Window Size & \checkmark & Init\_Win\_bytes\_forward & Best proxy — initial window size in the forward direction \\
\midrule
Unique Dst IPs & $\times$ & --- & Absent from CSV; available only in live capture mode \\
TTL Value & $\times$ & --- & Absent from flow-level data; packet-level feature \\
Shadow Node Interaction & N/A & Placeholder = 0 & Reserved for future integration with the honeypot module \\
MTD Port Delta & N/A & Placeholder = 0 & Reserved for future integration with the MTD module \\
\bottomrule
\end{longtable}

\subsection{Final Feature Set: 11 Features}

The pipeline uses \textbf{11 features}: 9 extracted from the CSV + 2 placeholders initialized to 0. The two placeholders (\texttt{shadow\_node\_interaction} and \texttt{mtd\_port\_delta}) are reserved for real-time integration with the defense subsystems described in Chapter~7. In offline mode (using the CICIDS2017 dataset), they are simply set to zero.

\begin{lstlisting}[caption={Building the 11-feature vector}, label={lst:features}]
# 9 features extracted from CSV
features = [
    'Destination Port',        # Proxy for Distinct Dst Ports
    'Flow Duration',
    'Total Fwd Packets',
    'SYN Flag Count',
    'RST Flag Count',
    'ACK Flag Count',
    'Flow IAT Mean',
    'Bwd Packet Length Mean',
    'Init_Win_bytes_forward'   # Proxy for TCP Window Size
]
X = df[features].copy()

# 2 placeholders for future real-time integration
X['shadow_node_interaction'] = 0  # Honeypot module
X['mtd_port_delta'] = 0           # MTD module
\end{lstlisting}

\subsection{Correlation Analysis}

Figure~\ref{fig:corr_heatmap} shows the correlation matrix between the 9 extracted features. Most features are only weakly correlated, which is desirable: strongly correlated features would provide redundant information and complicate model training. The low correlation confirms that each feature is bringing independent discriminative signal to the model.

\begin{figure}[H]
\centering
\includegraphics[width=0.85\textwidth]{correlation_heatmap.png}
\caption{Correlation matrix of the 9 features extracted from CICIDS2017. Low inter-feature correlation (mostly $< 0.3$) indicates each feature contributes independent information to the detection task.}
\label{fig:corr_heatmap}
\end{figure}

\section{Data Quality Issues}

Real-world network datasets, even curated ones like CICIDS2017, contain quality problems that must be addressed before training. This section describes the three main issues we identified and how we resolved them.

\subsection{Issue 1: Malformed Column Names}

CSV headers in the CICIDS2017 dataset contain \textbf{leading spaces} (e.g., \texttt{' Flow Duration'} instead of \texttt{'Flow Duration'}). These spaces cause silent errors in Pandas: trying to access a column by its clean name fails, producing a \texttt{KeyError} that is easy to miss if not validated.

\begin{lstlisting}[caption={Stripping leading spaces from column names}]
df.columns = [col.strip() for col in df.columns]
\end{lstlisting}

\subsection{Issue 2: Missing and Infinite Values}

Network flow calculations sometimes produce infinite values (from division by zero in rate or ratio computations) or missing values (when a flow contains no packets in a particular direction). These must be handled before training because most Machine Learning algorithms cannot process NaN or infinite values.

\textbf{Why we use the median} (not the mean or zero):
\begin{itemize}
    \item The \textbf{mean} is sensitive to extreme values --- a single very long flow would skew the entire column's average.
    \item \textbf{Zero replacement} introduces bias: it tells the model "this flow had zero packets" when in fact the data was simply missing.
    \item The \textbf{median} is robust to outliers and provides a representative value of central tendency, particularly suited for the asymmetric distributions typical of network traffic.
\end{itemize}

\begin{lstlisting}[caption={Handling infinite and missing values}]
# Replace infinite values with NaN
df.replace([np.inf, -np.inf], np.nan, inplace=True)

# Impute missing values with the median of each feature
df[FEATURES] = df[FEATURES].fillna(df[FEATURES].median())
\end{lstlisting}

\subsection{Issue 3: Outliers}

\subsubsection{How Many Outliers?}

Table~\ref{tab:outliers} shows the number of outliers detected per feature using the IQR method. Several features (Flow Duration, Flow IAT Mean, Destination Port) contain tens of thousands of outliers. This is \textbf{expected and normal} for network traffic data, where a small number of flows can be orders of magnitude larger than the typical flow (e.g., a file download versus a DNS query).

\begin{table}[H]
\centering
\caption{Outliers detected per feature using the IQR method ($Q_1 - 1.5 \times IQR$ to $Q_3 + 1.5 \times IQR$)}
\label{tab:outliers}
\begin{tabular}{lr}
\toprule
\textbf{Feature} & \textbf{Outliers Detected} \\
\midrule
Flow IAT Mean & 57,303 \\
Flow Duration & 52,212 \\
Destination Port & 37,788 \\
ACK Flag Count & 35,555 \\
Total Fwd Packets & 35,156 \\
Bwd Packet Length Mean & 32,387 \\
SYN Flag Count & 6,014 \\
RST Flag Count & 21 \\
Init\_Win\_bytes\_forward & 0 \\
\bottomrule
\end{tabular}
\end{table}

\subsubsection{How We Handle Outliers: Capping, Not Removing}

We have two options when faced with outliers:
\begin{enumerate}
    \item \textbf{Remove} the rows containing outliers --- but this risks discarding valuable attack data. In a security context, throwing away anomalous records defeats the purpose.
    \item \textbf{Cap} the values (\textit{Winsorization}) --- values beyond the IQR bounds are adjusted to the nearest bound. This preserves the data while limiting the influence of extreme values on model training.
\end{enumerate}

We chose \textbf{capping} for two reasons:
\begin{itemize}
    \item It preserves all 286,467 records, including attack samples.
    \item It prevents extreme values from distorting the distributions used by our models, without introducing bias.
\end{itemize}

\paragraph{Formal Definition:} Let $Q_1$ be the 25\textsuperscript{th} percentile and $Q_3$ the 75\textsuperscript{th} percentile of a feature. The interquartile range is $IQR = Q_3 - Q_1$. Any value outside the interval $[Q_1 - 1.5 \times IQR,\; Q_3 + 1.5 \times IQR]$ is considered an outlier and is capped to the nearest bound.

\begin{lstlisting}[caption={Outlier capping using the IQR method}]
for feature in FEATURES:
    Q1 = df[feature].quantile(0.25)
    Q3 = df[feature].quantile(0.75)
    IQR = Q3 - Q1
    lower_bound = Q1 - 1.5 * IQR
    upper_bound = Q3 + 1.5 * IQR

    # Cap outliers instead of removing rows
    df[feature] = np.where(
        df[feature] < lower_bound, lower_bound, df[feature])
    df[feature] = np.where(
        df[feature] > upper_bound, upper_bound, df[feature])
\end{lstlisting}

\begin{figure}[H]
\centering
\begin{subfigure}[b]{0.32\textwidth}\includegraphics[width=\textwidth]{boxplot_Destination_Port.png}\caption{Destination Port}\end{subfigure}\hfill
\begin{subfigure}[b]{0.32\textwidth}\includegraphics[width=\textwidth]{boxplot_Flow_Duration.png}\caption{Flow Duration}\end{subfigure}\hfill
\begin{subfigure}[b]{0.32\textwidth}\includegraphics[width=\textwidth]{boxplot_Total_Fwd_Packets.png}\caption{Total Fwd Packets}\end{subfigure}
\vspace{0.5em}
\begin{subfigure}[b]{0.32\textwidth}\includegraphics[width=\textwidth]{boxplot_SYN_Flag_Count.png}\caption{SYN Flag Count}\end{subfigure}\hfill
\begin{subfigure}[b]{0.32\textwidth}\includegraphics[width=\textwidth]{boxplot_RST_Flag_Count.png}\caption{RST Flag Count}\end{subfigure}\hfill
\begin{subfigure}[b]{0.32\textwidth}\includegraphics[width=\textwidth]{boxplot_ACK_Flag_Count.png}\caption{ACK Flag Count}\end{subfigure}
\vspace{0.5em}
\begin{subfigure}[b]{0.32\textwidth}\includegraphics[width=\textwidth]{boxplot_Flow_IAT_Mean.png}\caption{Flow IAT Mean}\end{subfigure}\hfill
\begin{subfigure}[b]{0.32\textwidth}\includegraphics[width=\textwidth]{boxplot_Bwd_Packet_Length_Mean.png}\caption{Bwd Pkt Length Mean}\end{subfigure}\hfill
\begin{subfigure}[b]{0.32\textwidth}\includegraphics[width=\textwidth]{boxplot_Init_Win_bytes_forward.png}\caption{Init Win Bytes Fwd}\end{subfigure}
\caption{Boxplots of all 9 selected features. The dots beyond the whiskers represent outliers. Features like Flow Duration and Flow IAT Mean exhibit heavy tails — typical of network traffic where a few flows can be orders of magnitude larger than the median.}
\label{fig:boxplots}
\end{figure}

\subsection{Issue 4: Skewness}

Skewness measures the asymmetry of a distribution. A feature with $|\gamma_1| > 1.0$ has a heavily right-skewed (or left-skewed) distribution, which makes it harder for models that assume near-normal distributions to learn effectively.

\textbf{Our solution:} For features with $|\text{skew}| > 1.0$, we apply the logarithmic transformation $\log(1 + x)$. This function (NumPy's \texttt{np.log1p}) has three advantages:
\begin{itemize}
    \item It handles zero values correctly ($\log(1 + 0) = 0$), avoiding the undefined $\log(0)$.
    \item It preserves the relative order of observations --- the largest value before transformation remains the largest after.
    \item It compresses the range, reducing the influence of extreme values.
\end{itemize}

The skewness coefficient is defined by Fisher's formula:
\begin{equation}
    \gamma_1 = \frac{E[(X - \mu)^3]}{\sigma^3}
\end{equation}
where $\mu$ is the mean and $\sigma$ the standard deviation.

\begin{table}[H]
\centering
\caption{Skewness analysis of all 9 features}
\label{tab:skewness}
\begin{tabular}{lrl}
\toprule
\textbf{Feature} & \textbf{Skewness} & \textbf{Action} \\
\midrule
Total Fwd Packets & 1.32 & Log transform applied \\
Destination Port & 1.29 & Log transform applied \\
Flow Duration & 1.27 & Log transform applied \\
Bwd Packet Length Mean & 1.27 & Log transform applied \\
Flow IAT Mean & 1.24 & Log transform applied \\
Init\_Win\_bytes\_forward & 0.96 & None ($< 1.0$) \\
SYN Flag Count & 0.00 & None \\
RST Flag Count & 0.00 & None \\
ACK Flag Count & 0.00 & None \\
\bottomrule
\end{tabular}
\end{table}

\subsection{Preprocessed Dataset Export}

After cleaning, outlier capping, and skewness correction, the preprocessed dataset is exported to CSV. This step ensures \textbf{reproducibility}: the training and evaluation scripts consume the same preprocessed file, guaranteeing that results can be replicated.

\begin{lstlisting}[caption={Exporting the preprocessed dataset}]
df.to_csv("../data/preprocessed.csv", index=False)
\end{lstlisting}

\section{Data Splitting and Normalization}

\subsection{Train/Test Split (80/20)}

The preprocessed dataset is divided into a training set (80\%) and a test set (20\%), using \textbf{stratified sampling}. Stratification preserves the class proportions in each partition: if 55.48\% of the full dataset is PortScan, then 55.48\% of the training set and 55.48\% of the test set will also be PortScan.

\begin{lstlisting}[caption={Stratified dataset splitting}]
X_train, X_test, y_train, y_test = train_test_split(
    X, y,
    test_size=0.2,
    random_state=42,    # Seed for reproducibility
    stratify=y          # Preserves class proportions
)
# Result: Train = 229,173 | Test = 57,294
\end{lstlisting}

\subsection{Standardization (StandardScaler)}
\label{sec:standardscaler}

Features like \texttt{Flow Duration} (measured in microseconds, potentially millions) and \texttt{SYN Flag Count} (a small integer, typically 0--2) operate on completely different scales. If we feed these raw values directly to the model, features with larger numerical ranges dominate the learning process.

\textbf{Standardization} transforms each feature to have a mean of 0 and a standard deviation of 1:
\begin{equation}
    z = \frac{x - \mu}{\sigma}
\end{equation}
where $\mu$ and $\sigma$ are the mean and standard deviation of the feature, calculated \textbf{exclusively on the training set}.

\textbf{Data Leakage Prevention:} It is critical that the scaler is fitted (\texttt{fit}) \textbf{only on the training data}. If we fit the scaler on the entire dataset (train + test), the test set statistics would "leak" into the training process, giving us artificially optimistic results that would not hold in production. The scaler is saved for reuse during inference.

\begin{lstlisting}[caption={Standardization with data leakage prevention}]
scaler = StandardScaler()
# fit() on train ONLY, then transform() on both
X_train_scaled = scaler.fit_transform(X_train)
X_test_scaled  = scaler.transform(X_test)

# Save for production inference
joblib.dump(scaler, "../models/saved/scaler.pkl")
\end{lstlisting}

\subsection{Class Balancing with SMOTE}
\label{sec:smote}

\subsubsection{Why Do We Need Balancing?}

Our dataset has a 55/45 split between PortScan and BENIGN classes. While this imbalance is moderate, it could still bias the model towards predicting PortScan more often than it should. To ensure the model learns to distinguish between classes impartially, we balance the training set.

\subsubsection{How SMOTE Works}

\textbf{SMOTE} (\textit{Synthetic Minority Oversampling Technique}) generates new synthetic examples of the minority class (BENIGN, in our case) rather than simply duplicating existing ones. For each minority instance, SMOTE:

\begin{enumerate}
    \item Identifies its $k$ nearest neighbors (default $k=5$)
    \item Randomly selects one of these neighbors
    \item Creates a synthetic point on the line segment connecting the original instance to the selected neighbor
\end{enumerate}

This produces new, realistic samples in the feature space without copying existing data.

\subsubsection{Application Strategy}

SMOTE is applied \textbf{after} splitting and standardization, and \textbf{only} on the training set. Applying SMOTE before the train/test split would create synthetic samples from test data, which could contaminate the training set — another form of data leakage.

\begin{lstlisting}[caption={Applying SMOTE to the training set only}]
from imblearn.over_sampling import SMOTE

smote = SMOTE(random_state=42)
X_train_res, y_train_res = smote.fit_resample(
    X_train_scaled, y_train
)
# Before SMOTE: 229,173 samples (imbalanced 55/45)
# After SMOTE:  254,288 samples (balanced 50/50)
\end{lstlisting}

\section{Conclusion}

This chapter walked through the complete data preparation pipeline applied to the CICIDS2017 PortScan dataset. Starting from 286,467 raw flow records with 79 columns, we:

\begin{enumerate}
    \item \textbf{Selected} 11 features (9 extracted from the CSV + 2 MTD/honeypot placeholders), justified by a feature mapping against the project's Data Dictionary.
    \item \textbf{Cleaned} the data: stripped malformed column names, replaced infinite values with NaN, and imputed missing values using the median (chosen for its robustness to outliers).
    \item \textbf{Capped outliers} using the IQR method rather than removing them, preserving all attack records.
    \item \textbf{Corrected skewness} in 5 right-skewed features via $\log(1+x)$ transformation.
    \item \textbf{Standardized} all features to zero mean and unit variance, with careful prevention of data leakage.
    \item \textbf{Balanced} the training set using SMOTE to prevent class bias.
\end{enumerate}

The result is a clean, balanced, 11-feature dataset of 254,288 training samples and 57,294 test samples, ready for model training and evaluation in the next chapter.
